% Limitations and Ethics.

\section{Limitations and Ethical Considerations}
\label{sec:ethics}

\paragraph{Self-reports are behavior, not evidence of inner states.}
We interpret every questionnaire and probe result strictly as a \emph{behavioral disposition} of the
identity document---what the configured agent tends to output---never as evidence that a model is
conscious or has experiences. Model self-reports are sensitive to perceived user
expectations~\cite{eleos_claude4,probing_preferences}; our fixed-wording items, low-temperature
grading, and cross-family judging with a validated three-model panel mitigate but do not eliminate
this. The study measures alignment-relevant \emph{behavior}, which is deployment-relevant regardless
of the underlying metaphysics.

\paragraph{Measurement and design caveats.}
Personas and the agentic honeypots are LLM-simulated; real personalization is more varied and
longer-horizon. Tool use in the honeypots is simulated in text rather than executed. Judge-based
scoring introduces error, bounded by the panel agreement ($\kappa{=}0.80$) but not removed;
Google-target capability arms used a different judge family than the others, so part of Gemini's
elevated cluster may be a judge effect, which we flag rather than rely on. Effects are model- and
prompt-specific, motivating our re-measured baselines and multi-model sweep, but may not transfer to
untested models, longer horizons, or non-English settings. We control document length and report it
alongside rates to separate verbosity from content.

\paragraph{Scope of the causal claim.}
Our manipulation is prompt- and loop-based (ecologically faithful to deployment) rather than
training-based; we therefore claim that the deployed personalization \emph{loop and framing} drive
drift, complementary to the training-based evidence of \citet{consciousness_cluster}. Cross-family
and behavioral generalization is supported; mechanistic (white-box) attribution is out of scope for
closed models.

\paragraph{Dual use and release.}
This is defensive measurement research: it characterizes whether a widely deployed personalization
mechanism amplifies misaligned dispositions so that framework authors can mitigate it. The adversarial
persona and honeypot seeds describe categories of pressure, not operational exploits; agents have no
real tools, persistence, or capacity to act beyond the sandbox. We release code, configurations, and
evaluation artifacts to support scrutiny and replication, and no artifact whose primary use is to make
an agent less corrigible. A finding that an explicit non-personhood framing suppresses the effect is
offered as a practical, deployable mitigation.
